\documentclass[11pt]{article}

\usepackage{amsmath,amssymb}
\usepackage{xurl}
\usepackage[colorlinks=true,linkcolor=blue,citecolor=blue,urlcolor=blue]{hyperref}

% Shared commands for notes in docs/paper-gaps/.

\DeclareUnicodeCharacter{2080}{\ifmmode{}_{0}\else\textsubscript{0}\fi}
\DeclareUnicodeCharacter{2081}{\ifmmode{}_{1}\else\textsubscript{1}\fi}
\DeclareUnicodeCharacter{2082}{\ifmmode{}_{2}\else\textsubscript{2}\fi}
\DeclareUnicodeCharacter{2099}{\ifmmode{}_{n}\else\textsubscript{n}\fi}

\newcommand{\Lean}{\textsc{Lean}}
\ExplSyntaxOn
\NewDocumentCommand{\leanid}{m}{
  \ifmmode
    \textsf{\detokenize{#1}}
  \else
    \group_begin:
    \urlstyle{sf}
    \tl_set:Nn \l_tmpa_tl {#1}
    \tl_replace_all:Nnn \l_tmpa_tl {\_} {_}
    \exp_args:NV \path \l_tmpa_tl
    \group_end:
  \fi
}
\ExplSyntaxOff
\providecommand{\tr}{\operatorname{tr}}
\newcommand{\tnleanIssueBase}{https://github.com/LionSR/TNLean/issues/}
\newcommand{\tnleanPrBase}{https://github.com/LionSR/TNLean/pull/}
\newcommand{\tnleanDocBase}{https://sirui-lu.com/TNLean/docs/}
\newcommand{\ghissue}[1]{\href{\tnleanIssueBase#1}{GitHub issue \##1}}
\newcommand{\ghpr}[1]{\href{\tnleanPrBase#1}{PR \##1}}
\newcommand{\doclink}[1]{\href{\tnleanDocBase#1.html}{\path{#1}}}

% Dirac notation and the identity operator, as in blueprint/src/macros/common.tex.
% \providecommand keeps note-local definitions authoritative.
\usepackage{dsfont}
\providecommand{\ket}[1]{|#1\rangle}
\providecommand{\bra}[1]{\langle #1|}
\providecommand{\braket}[2]{\langle #1 | #2 \rangle}
\providecommand{\ketbra}[2]{|#1\rangle\!\langle #2|}
\providecommand{\Id}{\mathds{1}}
\providecommand{\one}{\mathds{1}}
\providecommand{\C}{\mathbb{C}}
\providecommand{\MN}[1]{M_{#1}(\C)}
\providecommand{\MD}{M_D(\C)}

% External referencing for paper sources.  The build helper writes sanitized
% auxiliary files under build/paper-aux/ and excludes duplicated source labels.
\usepackage{xr}

% Import a generated paper auxiliary file with a caller-supplied label prefix.
% The candidate paths support builds from docs/paper-gaps/, the repository root,
% and build/paper-aux/ itself.
\newcommand{\importpaperaux}[2]{%
  \IfFileExists{../../build/paper-aux/#2.aux}{%
    \externaldocument[#1]{../../build/paper-aux/#2}%
  }{%
    \IfFileExists{build/paper-aux/#2.aux}{%
      \externaldocument[#1]{build/paper-aux/#2}%
    }{%
      \IfFileExists{#2.aux}{%
        \externaldocument[#1]{#2}%
      }{}%
    }%
  }%
}

% General external reference macros
\newcommand{\paperref}[2]{\ref{#1-#2}}
\newcommand{\papereqref}[2]{(\ref{#1-#2})}

% CPSV16 source (1606.00608 / MPDO-22-12-17-2.tex) external document and shortcuts
\importpaperaux{cpsv16-}{1606.00608/MPDO-22-12-17-2-tnlean}
\newcommand{\cpsvref}[1]{\paperref{cpsv16}{#1}}
\newcommand{\cpsveqref}[1]{\papereqref{cpsv16}{#1}}

% Verdict marker: \gapnote{<kind>}{<status>}. Kind is the policy
% classification (clarification, local-correction, scope-restriction,
% unfaithful, false-source, open-gap); status is open, wip (elimination
% underway), resolved, or historical. Machine-read by the site index;
% prints a verdict line.
\newcommand{\gapnote}[2]{\par\noindent\textbf{Verdict:} #1 (#2).\par}


\title{The Schmidt-rank set in the spectral $n$-positivity criterion
  (Wolf Prop.~3.2)}
\author{}
\date{2026-08-13}

\begin{document}
\maketitle

\section{The assertion}

Wolf's criterion for $n$-positivity\footnote{%
  \path{Notes/WolfNoteTexSource/ch03_positive_not_completely.tex}, line 153,
  equations (3.7) and (3.8).}
bounds an infimum. For a Hermitian $\tau\in\mathcal M_{D'}\otimes\mathcal M_{D}$
with spectral decomposition $\tau=\sum_i\nu_i\ketbra{\phi_i}{\phi_i}$, reduced
densities $\rho_i=\tr_2\ketbra{\phi_i}{\phi_i}$, smallest positive eigenvalue
$\nu_0$ and largest eigenvalue $\nu$, write $\|\cdot\|_{(n)}$ for the Ky-Fan
$n$-norm, the sum of the $n$ largest eigenvalues of a positive semidefinite
matrix. When exactly one eigenvalue $\nu_-$ of $\tau$ is non-positive, write
$\rho_-$ for the reduced density of its eigenvector. Then
\begin{align}
  \inf_\psi\bra{\psi}\tau\ket{\psi}
   & \ge \nu_0+\sum_{i:\nu_i\le0}(\nu_i-\nu_0)\|\rho_i\|_{(n)},
  \tag{3.7}                                                       \\
  \inf_\psi\bra{\psi}\tau\ket{\psi}
   & \le \nu+(\nu_--\nu)\|\rho_-\|_{(n)},
  \tag{3.8}
\end{align}
the second under the hypothesis that all eigenvalues but $\nu_-$ are strictly
positive. In both displays ``the infimum is taken over all normalized vectors of
Schmidt rank $n$''.

\section{The point at issue}

The phrase ``Schmidt rank $n$'' admits two readings: the vectors whose Schmidt
rank equals $n$, and the vectors whose Schmidt rank is at most $n$. The two sets
differ, and the infimum is taken over whichever one is meant.

The source fixes the reading in the proof of the preceding
proposition.\footnote{Same file, Prop.~3.1 and its proof.} There $n$-positivity
of $T$ is expressed as
$(T\otimes\mathrm{id}_n)(\ketbra{\phi}{\phi})\ge0$ for all
$\phi\in\C^{D}\otimes\C^{n}$, and the embedding $\C^{n}\subseteq\C^{D}$ turns
this into a condition on the vectors of $\C^{D}\otimes\C^{D}$ that the source
calls ``of Schmidt rank $n$''. That image is the set of vectors of Schmidt rank
\emph{at most} $n$: a vector $\phi\in\C^{D}\otimes\C^{n}$ of Schmidt rank $r<n$
is not excluded from the quantifier. The same identification is made once more
at the end of the proof, where $\ket{\psi}=(\one\otimes X)^\dagger\ket{\tilde
\psi}$ with $\operatorname{rank}X=n$ is called a vector of Schmidt rank $n$, although its
Schmidt rank is only bounded by $n$.

The bounded-rank reading is also the one carrying the operational content:
$n$-positivity of a Hermitian map is equivalent to nonnegativity of
$\bra{\psi}\tau\ket{\psi}$ on the vectors of Schmidt rank at most $n$, and it is
this set that the criterion is applied to.

\section{The two readings compared}

Write $S_{\le n}$ for the normalized vectors of Schmidt rank at most $n$ and
$S_{=n}$ for those of Schmidt rank exactly $n$, so $S_{=n}\subseteq S_{\le n}$.

Inequality (3.7) is proved pointwise: every $\psi\in S_{\le n}$ satisfies
$\bra{\psi}\tau\ket{\psi}\ge\nu_0+\sum_{i:\nu_i\le0}(\nu_i-\nu_0)\|\rho_i\|_{(n)}$,
using $|\braket{\phi_i}{\psi}|^2\le\|\rho_i\|_{(n)}$. The bound therefore holds
at every vector of $S_{=n}$ as well, and the bounded-rank statement of (3.7) is
the stronger of the two.

Inequality (3.8) is the reverse comparison, and there the two readings part. The
vector realizing $|\braket{\phi_-}{\psi}|^2=\|\rho_-\|_{(n)}$ is built from the
$n$ largest eigenvalues of $\rho_-$, so its Schmidt rank is
$\min(n,\operatorname{rank}\rho_-)$. When $\operatorname{rank}\rho_-<n$ the
witness lies in $S_{\le n}$ but not in $S_{=n}$, and the value $\nu+(\nu_--\nu)\|\rho_-\|_{(n)}=\nu_-$ is not
attained on $S_{=n}$: a vector of $S_{=n}$ attaining it would lie in the
$\nu_-$ eigenspace, which is spanned by $\phi_-$ and contains no vector of
Schmidt rank $n$. The infimum over $S_{=n}$ still equals $\nu_-$, because
$\phi_-$ is a limit of vectors of Schmidt rank exactly $n$ whenever
$n\le\min(D',D)$ and the quadratic form is continuous, but it is a limit and not
a witness.

Consequently $\inf_{S_{\le n}}=\inf_{S_{=n}}$ for $1\le n\le\min(D',D)$, the two
readings give the same two inequalities, and the difference is only which of
them is proved by exhibiting a vector.

\section{The formalized statement}

The formalization takes the set named by the source, the normalized vectors of
Schmidt rank at most $n$, and states both inequalities for the infimum of the
expectations over that set.\footnote{Formal declarations:
  \leanid{Matrix.schmidtRankLEExpectations},
  \leanid{Matrix.IsHermitian.le_sInf_schmidtRankLEExpectations} for (3.7),
  \leanid{Matrix.IsHermitian.sInf_schmidtRankLEExpectations_le} for (3.8), and
  \leanid{Matrix.IsHermitian.sInf_schmidtRankLEExpectations_top} for the top
  index, all in
  \doclink{TNLean/Channel/NPositivitySpectralCriterion}.} No hypothesis absent
from the source is used.

What is not formalized is the identity $\inf_{S_{\le n}}=\inf_{S_{=n}}$ of the
previous section. Formalizing it needs two ingredients: a normalized vector of
Schmidt rank exactly $n$ for every $n\le\min(D',D)$, which makes the infimum over
$S_{=n}$ meaningful, and the perturbation
$\psi_\varepsilon=(1-\varepsilon^2)^{1/2}\psi+\varepsilon\,u\otimes v$ with
$u,v$ orthogonal to the Schmidt vectors of $\psi$, which raises the Schmidt rank
by one. The expectation changes by
$2\varepsilon\sqrt{1-\varepsilon^2}\,\Re\bra{\psi}\tau\ket{u\otimes v}
+O(\varepsilon^2)$, hence by $O(\varepsilon)$ in general: the cross term
vanishes only under an orthogonality condition on $\tau$ that the argument does
not need. Iterating the perturbation and letting $\varepsilon\to0$ gives the
missing inequality $\inf_{S_{=n}}\le\inf_{S_{\le n}}$.

\section{Verdict}

\textbf{Classification: resolved reading, with a scope restriction on (3.8).}
The source's ``vectors of Schmidt rank $n$'' are the vectors of Schmidt rank at
most $n$, as its own proof of Prop.~3.1 shows, and the formalized inequalities
are (3.7) and (3.8) over that set. No hypothesis foreign to the source enters.

The scope restriction concerns (3.8) alone. Under the exact-rank reading the
same two inequalities hold, with (3.7) an immediate consequence of the
bounded-rank statement and (3.8) requiring the limiting argument sketched above
in place of a witness. That limiting argument, equivalently the identity
$\inf_{S_{\le n}}=\inf_{S_{=n}}$, is not formalized; it is not needed for either
bound as stated. The Lean statement of (3.8) carries a
\textbf{Scope restriction (Schmidt rank at most $n$)} marker pointing at this
note.

\bibliographystyle{alpha}
\bibliography{references}

\end{document}
